\documentclass[10pt,letterpaper]{article}
\usepackage[top=0.85in,left=2.75in,footskip=0.75in]{geometry}

% amsmath and amssymb packages, useful for mathematical formulas and symbols
\usepackage[right]{lineno}
\usepackage{amsmath,amssymb}

% Use adjustwidth environment to exceed column width (see example table in text)
\usepackage{changepage}

% cite package, to clean up citations in the main text. Do not remove.
\usepackage{cite}

% ligatures disabled
\usepackage[nopatch=eqnum]{microtype}
\DisableLigatures[f]{encoding = *, family = * }

% color can be used to apply background shading to table cells only
\usepackage[table]{xcolor}

% array package and thick rules for tables
\usepackage{array}

% create "+" rule type for thick vertical lines
\newcolumntype{+}{!{\vrule width 2pt}}

% create \thickcline for thick horizontal lines of variable length
\newlength\savedwidth
\newcommand\thickcline[1]{%
  \noalign{\global\savedwidth\arrayrulewidth\global\arrayrulewidth 2pt}%
  \cline{#1}%
  \noalign{\vskip\arrayrulewidth}%
  \noalign{\global\arrayrulewidth\savedwidth}%
}

% \thickhline command for thick horizontal lines that span the table
\newcommand\thickhline{\noalign{\global\savedwidth\arrayrulewidth\global\arrayrulewidth 2pt}%
\hline
\noalign{\global\arrayrulewidth\savedwidth}}

% Remove comment for double spacing
\usepackage{setspace} 
%\doublespacing

% Text layout
\raggedright
\setlength{\parindent}{0.5cm}
\textwidth 5.25in 
\textheight 8.75in

% Bold the 'Figure #' in the caption and separate it from the title/caption with a period
% Captions will be left justified
\usepackage[aboveskip=1pt,labelfont=bf,labelsep=period,justification=raggedright,singlelinecheck=off]{caption}
\renewcommand{\figurename}{Fig}

% Use the PLoS provided BiBTeX style
\bibliographystyle{plos2015}

% Remove brackets from numbering in List of References
\makeatletter
\renewcommand{\@biblabel}[1]{\quad#1.}
\makeatother

% Header and Footer with logo
\usepackage{lastpage,fancyhdr,graphicx}
\usepackage{epstopdf}
%\pagestyle{myheadings}
\pagestyle{fancy}
\fancyhf{}
%\setlength{\headheight}{27.023pt}
%\lhead{\includegraphics[width=2.0in]{PLOS-submission.eps}}
\rfoot{\thepage/\pageref{LastPage}}
\renewcommand{\headrulewidth}{0pt}
\renewcommand{\footrule}{\hrule height 2pt \vspace{2mm}}
\fancyheadoffset[L]{2.25in}
\fancyfootoffset[L]{2.25in}
\lfoot{\today}

% additional packages
\usepackage{etoolbox} %% <- for \cspreto, \csappto
\usepackage{amsfonts}
\usepackage{graphicx}
\usepackage{color}
\usepackage{url}
\usepackage{authblk}
\usepackage{bm}

% my commands
\DeclareMathOperator{\Var}{Var}
\DeclareMathOperator{\Cov}{Cov}
\newcommand{\aprcomment}[1]{{\textcolor{blue}{APR: #1}}}
\newcommand{\E}{\mathbb{E}}
\usepackage{xspace}
\newcommand{\moments}{\texttt{moments}\xspace}
\newcommand{\fwdpy}{\texttt{fwdpy11}\xspace}
\newcommand{\tskit}{\texttt{tskit}\xspace}

%% Patch 'normal' math environments:
\newcommand*\linenomathpatch[1]{%
  \cspreto{#1}{\linenomath}%
  \cspreto{#1*}{\linenomath}%
  \csappto{end#1}{\endlinenomath}%
  \csappto{end#1*}{\endlinenomath}%
}

\linenomathpatch{equation}
\linenomathpatch{gather}
\linenomathpatch{multline}
\linenomathpatch{align}
\linenomathpatch{alignat}
\linenomathpatch{flalign}

\linenumbers

\title{Archaic introgression and the distribution
of shared variation under stabilizing selection}
\author[]{Aaron P. Ragsdale}
\affil[]{apragsdale@wisc.edu}
\affil[]{Department of Integrative Biology, University of Wisconsin--Madison, WI 53706}
\date{\today}

\begin{document}
\maketitle

\begin{abstract}

    \noindent Many phenotypic traits are under stabilizing selection, which
    maintains a population's mean phenotypic value near some optimum. The
    dynamics of traits and trait architectures under stabilizing selection have
    been extensively studied for single populations at steady state. However,
    natural populations are seldom at steady state and are often structured in
    some way. Admixture and introgression events may be common, including over
    human evolutionary history. Because stabilizing selection results in
    selection against the minor allele at a trait-affecting locus, alleles from
    the minor parental ancestry will be selected against after admixture. We
    show that the site-frequency spectrum can be used to model the genetic
    architecture of such traits, allowing for the study of trait architecture
    dynamics in complex multi-population settings. We use a simple
    deterministic two-locus model to predict the reduction of introgressed
    ancestry around trait-contributing loci. From this and individual-based
    simulations, we show that introgressed-ancestry is depleted around such
    loci. When introgression between two diverged populations occurs in both
    directions, as has been inferred between humans and Neanderthals, the
    locations of such regions with depleted introgressed ancestry will tend to
    be shared across populations. We argue that stabilizing selection for
    shared phenotypic optima may explain recent observations in which regions
    of depleted human-introgressed ancestry in the Neanderthal genome overlap
    with Neanderthal-ancestry deserts in humans.

\end{abstract}

\begin{center}
    \subsubsection*{Author Summary}
\end{center}
\small
\begin{quote}

    After admixture between two diverged populations, selection often removes
    introgressed haplotypes. Selection against introgressed ancestry has
    multiple proposed causes, including accumulated deleterious variation in
    the source population, locally adapted variation being more fit, or genomic
    incompatabilities. In this article, we find that normal selection on
    complex traits, that is, stabilizing selection, causes selection against
    introgressed ancestry. We develop theory that enables predictions for
    patterns of introgressed ancestry under stabilizing selection on
    phenotypes. Using this, we hypothesize that the distribution of
    introgressed-ancestry ``deserts'' should be correlated between populations
    when there is bi-directional gene flow, as inferred between Neanderthals
    and humans.

\end{quote}

\normalsize

\onehalfspacing

\section*{Introduction}

Genomic surveys of natural systems show that historical admixture among
diverged populations and closely related taxa commonly occurs
\cite{brandvain2014speciation, skoglund2015ancient, suvorov2022widespread} and
is widespread in primate \cite{tung2017contribution, sorensen2023genome} and
hominin \cite{wolf2018outstanding, peter2020100} evolution. Admixture is
therefore a frequent driver of phenotypic and molecular variation and can
contribute to the genetic architectures of complex traits. In humans, archaic
introgression from Neanderthals and Denisovans has attracted considerable
attention, including efforts to describe the historical processes leading to
observed distributions of introgressed DNA in present-day populations
\cite{prufer2014complete, villanea2019multiple, chen2020identifying} and the
contribution of introgressed variation to quantitative traits
\cite{sankararaman2016combined, wei2023lingering}.

Once introduced through admixture, introgressed alleles may be selected for or
against. Some introgressed haplotypes were likely positively selected in modern
\emph{Homo sapiens} (here, ``humans'') \cite{huerta2014altitude,
racimo2017signatures, enard2018evidence, gower2021detecting}, possibly due to
locally adaptive variation that provided fitness advantages as humans
encountered novel environments, and introgressed alleles may be maintained at
intermediate frequencies due to heterosis \cite{harris2016genetic}. Despite
some documented cases of adaptive introgression, most introduced functional
alleles likely were selected against in humans \cite{harris2016genetic,
juric2016strength, veller2023recombination}.

Since Neanderthal and Denisovan population sizes were relatively small for
hundreds of thousands of years, theory predicts they would have accumulated
deleterious variation at an increased rate \cite{kimura1963mutation,
lynch1990mutation}. Introgressed haplotypes loaded with more deleterious
mutations would have been rapidly removed by selection after admixture. Mapping
the distribution of Neanderthal-introgressed haplotypes in humans shows a
reduction of Neanderthal-related ancestry in coding and regulatory regions
\cite{petr2019limits, telis2020selection, yermakovich2023long}. These
``deserts'' of Neanderthal ancestry support the hypothesis that introgressed
functional alleles were selected against \cite{sankararaman2014genomic,
vernot2014resurrecting, sankararaman2016combined, vernot2016excavating}. The
term ``deserts'' in the literature has been used to describe both a finite
number of very large depletions of archaic ancestry extending many Mb in length
\cite{sankararaman2014genomic, vernot2014resurrecting,
sankararaman2016combined, vernot2016excavating} or depletions across smaller
regions such as protein-coding or non-coding regulatory loci
\cite{petr2019limits, telis2020selection, yermakovich2023long,
harris2023diverse}. Here, we use this term to refer to any region of depleted
ancestry from an introgressing population.

There is growing genetic evidence that Neanderthals reciprocally received
genetic material from early humans \cite{kuhlwilm2016ancient,
hubisz2020mapping, harris2023diverse, li2024recurrent}. This gene flow occurred
tens to hundreds of thousands of years prior to Neanderthal introgression in
humans during the global dispersal of modern humans around 60 ka. The earlier
admixture is supported by \emph{H. sapiens} outside of African around 200--100
ka \cite{schwarcz1988esr, grun2005u, harvati2019apidima, beyer2021climatic},
potentially overlapping with Neanderthals and providing opportunities for early
contacts. While estimates of the genomic contribution of early humans to
Neanderthals vary, around $6\%$ of later Neanderthal genomes may trace through
this admixture event \cite{harris2023diverse}. Under a ``load'' model, if
human-related haplotypes carried fewer deleterious alleles due to their larger
long-term effective population size, human-introgressed DNA would have been
favored in Neanderthal genomes. The replacement of Neanderthal mitochondrial
and Y chromosomes by early human haplotypes appears to support this model of
post-admixture positive selection in the Neanderthal lineage
\cite{posth2017deeply, petr2020evolutionary}.

Models for selection against introgressed alleles are often based on
deleterious load or hybrid incompatibilities \cite{muller1942isolating}. These
models, founded in population genetics theory, rarely take into account
selection operating on phenotypic traits or the relationship between genetic
and phenotypic variation.  Many phenotypic traits are thought to be under
stabilizing selection \cite{sanjak2018evidence, sella2019thinking}, including
gene regulation \cite{gilad2006natural, hodgins2015gene, price2022detecting}.
Because some of the strongest signals of selection against
Neanderthal-introduced alleles are in regulatory regions
\cite{sankararaman2014genomic}, stabilizing selection on quantitative traits
may be particularly relevant to the dynamics of functional genetic variation
after introgression among hominins.

Stabilizing selection acts to maintain the phenotypic distribution of a trait
near some optimum, which is achieved by reducing phenotypic variation
(Figure~\ref{fig:stab-sel}A). When the mean phenotype of the population is
close to the phenotypic optimum, classical models predict that the minor allele
at a trait-affecting locus is selected against, with allele-frequency dynamics
equivalent to underdominant selection \cite{robertson1956effect}. This has
proven to be a useful model for understanding genetic architectures of traits
under stabilizing selection in single-population settings (e.g.,
\cite{keightley1988quantitative, simons2018population, hayward2022polygenic}).

\begin{figure}[tb!]
    \centering
    \caption{
        \textbf{Additive genetic variance under stabilizing selection.}
        (A) Stabilizing selection acts to maintain phenotypic values of
        individuals in the population near some optimum. Throughout, we assume
        a Gaussian fitness function.
        (B) With low mutational variance ($V_M$), the expected additive
        genetic variance ($V_G$) is proportional to the
        population-scaled mutation rate. When $V_M$ is large ($N_e V_M > 1$) so that
        mutational effects can be strong, $V_G$ is independent of $V_M$.
        The stochastic house-of-cards model (Eq.~\ref{eq:SHC})
        interpolates these regimes, assuming steady-state dynamics
        \cite{burger1989much}. Expected $V_G$ computed using an SFS approach
        (developed here using \moments \cite{jouganous2017inferring})
        matches simulations assuming linkage equilibrium between loci affecting the
        trait, which align closely with the stochastic house-of-cards
        approximation.
        Here $V_S=1$, $N_e=10{,}000$, the mutation rate
        $\mu=0.01$ (per haploid), and mutation effects are drawn from a normal
        distribution with mean zero and variance $V_M$.
    }
    \label{fig:stab-sel}
\end{figure}

In diverged populations, the genetic variation contributing to a trait under
stabilizing selection has a higher rate of turnover compared to neutral
evolution \cite{yair2022population}. We therefore expect a rapid divergence of
trait architectures in the human and Neanderthal lineages at trait-contributing
loci, even when the mean phenotype in each population remains close to the same
trait optimum. When a derived allele at high frequency in one population is
introduced to another population in which it was previously absent, it will be
at low frequency (if the admixture proportion is low) and subsequently selected
against. Likewise, if the ancestral allele is reintroduced to a population
fixed for a derived allele, the ancestral allele will be at low frequency and
will be selected against. As we show below, selection acts against the
introgressed allele in either case, whether it is ancestral or derived and
regardless of the historical relative sizes of the populations involved. In
concurrent work, Veller and Simons (2024) \cite{veller2024stabilizing}
demonstrate this effect by deriving the expected decay of minor parental
ancestry under stabilizing selection after admixture. This prediction
contrasts with the population-genetics load model, in which haplotypes with
fewer deleterious variants (such as those from the population with larger
historical size) are favored after introgression in either direction
\cite{castellano2014patterns,petr2019limits,petr2020evolutionary}.

In this article, we show that admixture between diverged populations results in
selection against the minor parental ancestry at loci contributing to a trait
under stabilizing selection. We develop a numerical approach based on the
site-frequency spectrum to predict genetic variation under complex demographic
scenarios, which we use to partition predicted trait heritability by
introgressed and non-introgressed variation. Using simulations with linkage, we
demonstrate that deserts of introgressed ancestry form around
trait-contributing loci. When gene flow occurs bidirectionally, such deserts
will tend to overlap in location across populations. We argue that stabilizing
selection on shared trait optima may explain the overlap of
introgressed-ancestry deserts in human and Neanderthal genomes after reciprocal
introgression \cite{harris2023diverse}.

%\section*{Model and Methods}

\subsection*{Model and Theory}

We consider a polygenic trait for which an individual's additive genetic value
is the sum over all effects of alleles in their genome. For individual $i$,
\(G_i = \sum_l g_{l, i} a_l\), where $a_l$ is the effect size of the derived
allele at locus $l$, and \(g_{l,i}\in\{0,1,2\}\) is their genotype at that
locus (i.e., the number of derived alleles they carry). With linkage
equilibrium between trait-affecting loci, the expected additive genetic
variance is \(V_A = \sum_l 2p_l(1-p_l)a_l^2\), where $p_l$ is the allele
frequency at locus $l$. We ignore dominance and epistasis (often argued to be a
reasonable modeling choice, e.g., \cite{barton2002understanding,
zhang2002pleiotropic}), so that the genetic variance \(V_G = V_A\). We further
ignore environmental effects and noise; under the assumption of additive
environmental effects, this can be absorbed into $V_S$
\cite{lande1975maintenance, simons2018population}. Thus, only genetic effects
are considered, so that the phenotypic variance \(V_P=V_G\).

Stabilizing selection acts to reduce phenotypic variation around the optimum
value $O$, typically set to zero, and we assume a Gaussian fitness function
(with variance $V_S$, Figure~\ref{fig:stab-sel}A) so that relative fitness
(relative to an individual with optimal phenotype) is given by \(f(G_i | O,
V_S) = \exp{(-(G_i - O)^2 / 2 V_S)}\). $\frac{1}{V_S}$ is interpreted as the
strength of selection on the trait, so larger $V_S$ implies weaker selection.
For a population with mean phenotype at or very close to the optimum, the
mean fitness of the population [assuming a normal distribution of phenotypic
values in the population, \cite{turelli1994genetic, urban2013asymmetric}] is
\[\bar{w} \approx \int_{-\infty}^\infty
f(G | O, V_S) \mathcal{N}(O, V_G) dG =
\left(\frac{V_S}{V_S+V_G}\right)^{1/2},\]
so that as the genetic variance increases, mean fitness decreases (see
SI, Section~S1).

\subsubsection*{Mutation rates, effect sizes and genetic variance}

If all alleles contribute equally to the trait
with effect sizes \(\pm a\) occurring in equal proportion,
Keightley and Hill \cite{keightley1988quantitative} showed that the
dynamics of $V_G$ can be
approximated with the recursion
\[V_{G,t+1} \approx V_{G,t}\left(1-a^2/2(V_S+V_{G,t})\right)
\left(1-1/2N_e\right) + 2 \mu a^2,\]
where $\mu$ is the per-haploid, per-generation rate of new mutations
contributing the trait. In the
large-population-size limit, this gives the well-known result for
steady-state additive genetic variance,
\[V_G \approx 4\mu V_S,\] provided \(V_G \ll V_S\).

Mutations will not generally all have the same effect size, but rather are
drawn from some distribution. Here, when modeling mutation effects as
non-constant, we assume effect sizes follow a normal distribution with mean 0
and variance $V_M$. When population sizes or effect sizes are small, drift
dominates the dynamics of $V_G$, and at steady state Lande \cite{lande1976natural}
found \[V_G \approx 4 N_e \mu V_M.\] Interpolating between the drift- and
selection-dominated regimes,
\begin{align}\label{eq:SHC}
    V_G & \approx \frac{4 \mu V_S}{1 + \frac{V_S}{N_e V_M}}.
\end{align}
This, the ``stochastic
house-of-Cards'' (SHC) approximation (Figure~\ref{fig:stab-sel}B), was given by
Burger \cite{burger1989much} and is discussed in detail in Chapter 28 of
Walsh and Lynch \cite{walsh2018evolution}.

\subsubsection*{Approximating allelic dynamics via underdominance}

As initially shown by Robertson \cite{robertson1956effect} (see also,
\cite{keightley1988quantitative, simons2018population}), when the mean
phenotype of the population is close to the optimal phenotypic value,
stabilizing selection results in selection against the minor alleles at loci
contributing additively to the trait, with dynamics mirroring symmetric
underdominance. In general, for an allele at frequency $p$ with selection
coefficient $s$, the expected change in $p$ over one generation is
\begin{equation}
    \E[\Delta p] = sp(1-p).
\end{equation}

In our case, the selection coefficient $s$ depends on the strength of selection
on the trait $V_S$ and the effect size $a$, as well as the frequency of the
allele, so that
\begin{equation}\label{eq:selcoeff}
    s\approx -\frac{a^2(1-2p)}{2(V_S+V_G)} 
    = \frac{\left(p-\frac{1}{2}\right)a^2}{V_S + V_G}
    \approx \frac{\left(p-\frac{1}{2}\right)a^2}{V_S},
\end{equation}
when $V_G \ll V_S$. This model assumes linkage equilibrium between
trait-contributing loci. We note that LD is expected to be generated between
loci even in the absence of linkage \cite{bulmer1971effect}, and
Negm and Veller \cite{negm2024effect} recently showed how to
correct for its effect on the selection dynamics of a focal allele.

Because $a^2/(V_S+V_G)$ is always positive for any $a\not=0$, we see that
selection pushes allele frequencies to zero if $p<1/2$ and to one if $p>1/2$,
resulting in symmetric underdominance. Details are shown in the Supporting
Information, Section~S1.

\begin{figure}[tb!]
    \centering
    \caption{
        \textbf{Population size changes and additive genetic variance.}
        We used \moments to track expected genetic variance of a trait
        under stabilizing selection as the population changes size. (A) The
        population goes through a 10-fold size reduction, followed by a recovery.
        (B) All mutations have effect sizes $\pm0.01$, $0.04$, or $0.1$, with equal
        probability of being trait increasing or decreasing. Depending on the
        selection coefficient (Equation~\ref{eq:selcoeff}) compared to $N_e$, $V_G$
        could increase or decrease after a sudden population size change.
        Predictions using \moments match simulations with linkage equilibrium
        between trait-affecting loci
        (Figures~\ref{fig:one-pop-0.01}--\ref{fig:one-pop-0.1}).
    }
    \label{fig:one-pop}
\end{figure}

\begin{figure}[tb!]
    \centering
    \caption{
        \textbf{Genetic variance of a trait under stabilizing selection with
        multi-population demography.}
        (A,B) Two simple demographic models, in which Deme 1 has size 10,000
        and Deme 2 has size 1,000. In each scenario, we allow $5\%$ admixture
        from one deme to the other after begin isolated for 20,000 ($2N_e$)
        generations.
        (C,D) We compared predicted (additive) $V_G$ from the site-frequency
        spectra (using \moments) to simulations
        without linkage between trait-affecting loci. In Deme 2, $V_G$ decreases
        due to their population size reduction. After admixture in both
        directions, $V_G$ increases in the recipient deme
        and then decays to steady-state levels. Predictions from \moments closely
        match observed $V_G$ in simulations.
        Here, mutations were drawn from a normal distribution with $V_M=0.0025$.
        Other parameters: $\mu=0.025$, $V_S=1$, and the optimal phenotype remained
        the same in each population.
        (E,F) Using \moments, we partitioned additive contributions to $V_G$ by
        mutations that were previously segregating in the focal population,
        introgressed variants (either derived or reintroduced ancestral alleles),
        and by new mutations since the time of admixture. In both demographic
        scenarios, introgressed variants contribute a substantial proportion to
        $V_G$, though it is primarily composed of introduced derived alleles when
        admixture is from the small to the large population, while primarily
        reintroduced ancestral alleles under the reverse direction of gene flow.
        In both cases, $V_G$ is increasingly due to new mutations as the genetic
        architecture of the trait turns over with time since admixture.
    }
    \label{fig:toy-admixture}
\end{figure}

\section*{Results}


\subsection*{Additive genetic variance after admixture}

We typically expect genetic variance to increase after introgression. The
amount that genetic variance increases depends on the allelic differences
accumulated between populations and the effects of those alleles. Assuming
linkage equilibrium, and ignoring dominance and epistasis, \(V_G=\sum_l
2p_l(1-p_l)a_l^2\). After admixture, with proportion $f$ contributed by the
population labeled 0 and $1-f$ by population 1, \(p_l = fp_{l,0} +
(1-f)p_{l,1}\). Plugging into the expression for $V_G$ and after a bit of
algebra (SI, Section~S3), we can write the expected genetic
variance directly after admixture as
\begin{align}
V_G = fV_{G,0} + (1-f) V_{G,1} + 2f(1-f)\sum_l F_{2,l} a_l^2,
\end{align}
where $F_2 = (p_0 - p_1)^2$ is the squared difference in allele frequencies at
a locus \cite{peter2016admixture}. This result is known
(e.g., \cite{tufto2000quantitative}), showing that additive variance is equal
to that in the source populations weighted by their contributions, plus a term
that depends on the divergence at trait-affecting loci between the populations
weighted by the quadratic factor $2f(1-f)$. We note that $V_G$ is expected to
increase only in the second generation after admixture, so that selection
against introgressed alleles is not immediate \cite{veller2024stabilizing}.
This effect is not captured by this expression.

$F_2$ at a given locus depends on the demographic history
relating the two populations and the effect size at the locus due to selection
on the trait. In the infinitesimal limit, involving many loci each of
vanishingly small effect, dynamics at a given locus will be approximately
neutral, so that $F_2$ depends only on demographic history. In this case,
\begin{align} \label{eq:VG-admixture}
    V_G & \approx f V_{G,0} + (1-f) V_{G,1} + 2f(1-f)\mathbb{E}[F_2] \sum_l a_l^2 \\
    \nonumber
    & = f V_{G,0} + (1-f) V_{G,1} + 2f(1-f)\mathbb{E}[F_2] L V_M,
\end{align}
where $L$ is the number of trait affecting loci.


\begin{figure}[ht!]
    \centering
    \caption{
        \textbf{Allele contributions to heritability under human-Neanderthal
        reciprocal introgression.}
        (A) A model of divergence and admixture between humans and
        Neanderthals. Using \moments, we computed predicted $V_G$ over time,
        partitioned by variation that was introgressed vs. non-introgressed
        (Figure~\ref{fig:human-neand-VG}).
        (B-E) Predicted per-SNP contributions to genetic variance ($h^2$ per
        SNP) is plotted over the 50 thousand years following introgression.
        For non-introgressed variants, we also plot $h^2$-per-SNP weighted by
        allele frequencies matching those of introgressed variants. These
        are shown for (B,C) Neanderthal-to-human introgression 50 kya,
        (D,E) human-to-Neanderthal introgression 250 kya, (B,D) $V_M=0.0025$,
        and (C,E) $V_M=0.0001$.
    }
    \label{fig:human-neand-h2}
\end{figure}

\subsubsection*{Predicted \bm{$V_G$} from the SFS}

Expected allele frequency differences ($F_2$) for selected alleles differ from
neutrality. For negative and underdominant selection, $F_2$ is reduced relative
to neutral expectations (Figure~\ref{fig:underdominance-F2}). Because analytic
solutions are unavailable for arbitrary evolutionary scenarios involving
multiple populations, we numerically solve for the expected joint distribution
of allele frequencies (the SFS) before and after admixture. This provides a
numerical solution for expected $V_G$, which can be tracked over time
(Methods). Comparing to simulations with free recombination between loci, we
observe close agreement with average $V_G$ at all times
(Figure~\ref{fig:toy-admixture}C,D). In the scenarios tested here, admixture
causes a sudden increase in $V_G$ followed by a fairly rapid return to
pre-admixture levels, which is recovered by our numerical approach.

Modeling the dynamics of genetic variance using the SFS lets us examine
contributions to $V_G$ from different classes of mutations, such as those at
different frequencies or arising at different times, and how those
contributions change over time. In particular, we may quantify the contribution
to $V_G$ from alleles that were already segregating in the recipient
population, those that were introduced through introgression, and new mutations
since the time of admixture (Figure~\ref{fig:toy-admixture}E,F). In many
scenarios of interest, in which populations are considerably diverged at the
time of admixture, genetic architectures will be largely unique in each
population. After mixing, previously segregating and introgressed alleles each
contribute to $V_G$ before going to fixation or loss, and the variance of the
trait is increasingly due to new mutations. 

Introgressed variation can initially make up a considerable portion of $V_G$,
with those alleles being either newly introduced derived alleles or
reintroduced ancestral alleles. The relative sizes of the two populations
impact the numbers of each, as derived alleles will accumulate more readily in
a population with smaller effective size. Nonetheless, the overall increase in
$V_G$ is similar in both directions of introgression, as the symmetric term
\(2f(1-f)\sum_l F_{2,l} a_l^2\) contributes in either case and can be much
larger than $f V_G$ from the source population
(Equation~\ref{eq:VG-admixture}).

\subsubsection*{Complex demography and partitioning heritability by origin of
alleles}

We used a historical model resembling inferred human-Neanderthal history
(Figure~\ref{fig:human-neand-h2}A) to explore the effects of population size
changes and reciprocal admixture on the additive genetic architecture of traits
under stabilizing selection. As expected (Figure~\ref{fig:toy-admixture}),
population contractions decrease $V_G$ as drift removes allelic diversity at
trait-affecting loci, while introgression increases $V_G$.

Because the genetic architectures considered here are purely additive, we can
track mutations in an admixed populations by whether they were previously
segregating, fixed or lost in either parental populations or if they arose as
new mutations since the time of admixture. Partitioning $V_G$ by contributions
from these different sets of mutations (Figure~\ref{fig:human-neand-VG}), we
find it is still primarily contributed by previously segregating,
non-introgressed mutations. $V_G$ due to existing mutations decays
monotonically over time and is rapidly replaced by new mutations.

The average contribution of introgressed vs. non-introgressed SNPs to $V_G$
(i.e., $h^2$-per-SNP) can similarly be tracked over time. For the
human-Neanderthal demographic model and genetic architectures considered
here, the contribution per-SNP of introgressed variants is initially lower than
that of non-introgressed variants. These contributions change over time,
depending on mutational variance, as well as demography
(Figure~\ref{fig:human-neand-h2}). When weighting $h^2$-per-SNP of
non-introgressed SNPs by matching to allele frequencies of introgressed
variants, relative contributions depend sensitively on evolutionary parameters
and the time since admixture.

\subsection*{The effects of linkage}

In the preceding sections, we found that approximating the dynamics of
trait-affecting alleles using an underdominant selection model
\cite{robertson1956effect} provides an excellent approximation of $V_G$ in
complex demographic scenarios. However, this relies on linkage equilibrium
between trait-affecting alleles. The inclusion of linkage can cause noticeable
distortions of expected $V_G$, so that Equation~\ref{eq:SHC} differs from
observed $V_G$ at steady state \cite{burger1989much, burger1994distribution,
walsh2018evolution}.

To investigate the effects of linkage, we used chromosome-scale
individual-based simulations \cite{thornton2019polygenic}. By varying the
mutation rate and the variance of effect sizes of new mutations, we included
scenarios ranging from low to high polygenicity and from weak to strong
selection on individual alleles. In this and the following sections, we
highlight two effects of linkage. First, we observe deviations of $V_G$ from
expectations under free recombination, which can be large for highly polygenic
traits. Second, selection on introgressed trait-affecting alleles results in a
reduction of introgressed ancestry in surrounding regions.

With linkage between two or more loci contributing to a trait under stabilizing
selection, linkage disequilibrium (LD) can develop between alleles. Notably,
stabilizing selection will lead to coupling LD between mutations of
opposite-signed effects and repulsion LD between those of same-signed effects
\cite{bulmer1971effect}. This is expected to decrease $V_G$. For lower
mutational variances ($2N_e V_M\approx1$, Figure~\ref{fig:supp-low-VM}), we
observe such a reduction in $V_G$. With low mutational input, and thus low
polygenicity, $V_G$ in individual-based simulations closely matches
expectations under linkage equilibrium. With increasing mutation rate, $V_G$ is
reduced relative to those expectations. However, when the mutational variance
is much larger ($2N_eV_M\approx 50$), we see the reverse trend
(Figure~\ref{fig:supp-high-VM}). At low mutation rate, there is a close match
between observed $V_G$ and expectations, although simulated values are slightly
higher. As the mutation rate increases, $V_G$ increases to be relatively much
larger than expectations, rather than smaller.

The strength of the deviation of $V_G$ between models with and without linkage
depends on a number of factors. The total mutation rate affects not only the
polygenicity, with higher mutation rates producing more segregating alleles,
but also influences the mean recombination rate between alleles, as they will
be more densely distributed in the genome. The distribution of effect sizes
plays an important role, as seen by the opposite trends in $V_G$ distortions
(Figure~\ref{fig:supp-low-VM} vs. Figure~\ref{fig:supp-high-VM}), which are
most apparent with high mutation rates and thus high polygenicity. In these
cases, there appear to be complex dynamics involving the reduction in $V_G$ due
to the Bulmer Effect \cite{bulmer1971effect} and an increase in $V_G$ due to
selective interference \cite{hill1966effect}.

\begin{figure}[ht!]
    \centering
    \caption{
        \textbf{Reduced introgressed ancestry around alleles contributing to a
        trait under stabilizing selection.} Using a deterministic model
        (Equations~\ref{eq:system-p}--\ref{eq:system-D}), we model the
        frequencies of an introgressed trait-affecting allele and a linked
        neutral allele, initially absent from the recipient population so that
        their frequencies equal the introgression proportion ($f=0.05$). The
        neutral allele frequency is reduced at a rate that depends on the
        effect size of the selected allele and the probability of recombination
        between them. LD (as measured by $D=\Cov(p,q)$) decays to zero over
        time. (C) Alleles with strong effects are expected to result in a
        larger depletion of introgressed ancestry around the selected locus.
        (D,E) Compared to individual-based simulations, the deterministic model
        predicts the dip in introgressed ancestry around trait-affecting loci,
        when the mutation rate is low, so that trait-affecting loci are
        sparsely distributed. When mutation rates (and thus polygenicity) are
        high, selected alleles tend to be more tightly linked, so that
        selective interference is more pronounced and local ancestry is
        affected by multiple selected alleles.
    }
    \label{fig:linkage}
\end{figure}

\subsubsection*{Introgressed ancestry is reduced around introduced
trait-affecting alleles}

Because introgressed trait-affecting alleles are selected against as the minor
allele, introgressed ancestry segments in the regions surrounding the selected
alleles will also be removed due to linkage. The expected reduction in
introgressed ancestry will depend on the effect size of the linked
trait-affecting allele and the local probability of recombination. We first
consider a deterministic model for the frequency dynamics of introgressed
alleles (one selected, one neutral) with variable rate of recombination. This
simple model ignores the effects of drift and of interference between multiple
selected alleles.

We model a trait-affecting locus with a fixed difference between the two
parental populations, in which the derived allele may be fixed in either
population. At the functional locus, with admixture proportion $f$ from the
minor parental source, the initial frequency of the derived allele is either
$p_0=f$ or $p_0=1-f$. Over one generation, the expected allele frequency at the
selected locus is, to leading order in $s$,
\begin{align}\label{eq:system-p}
    p_{t+1} & = p_t + s_tp_t(1-p_t),
\end{align}
with \(s_t=(p_t-1/2)a^2/V_S\) for stabilizing selection
(Equation~\ref{eq:selcoeff}).

We consider a neutral locus separated from the selected locus by recombination
fraction $r$. Initially, the expected frequency of linked introgressed ancestry
is \(q_0=f\), which changes over time due to linked selection on the
trait-affecting allele. Letting \(D=\Cov(p,q)\) be the standard covariance
measure of LD between the alleles at the two loci, $q$ is expected to change as
\begin{align}\label{eq:system-q}
    q_{t+1} & = q_t + s_t D_t.
\end{align}
Initially, \(D_0=\pm f(1-f)\), with \(D\) being positive if \(p_0=f\) and
negative if \(p_0=1-f\). LD between the loci changes deterministically over
time due to both selection and recombination, so that
\begin{align}\label{eq:system-D}
    D_{t+1} = D_t - r D_t + s_t D_t (1-2p_t).
\end{align}
Together, this forms a simple nonlinear system of equations for the
deterministic change in allele frequencies at the two loci and LD between them.

Using this model, we predicted the changes in introgressed allele frequencies
and LD after admixture for given effect size $a$ and recombination rate $r$
(Figure~\ref{fig:linkage}A,B). As expected, smaller effect sizes result in a
slower decay in introgressed ancestry frequency at both selected and linked
loci, and larger recombination rates more quickly decouple the linked ancestry
from the selected allele dynamics. Thus, the expected reduction in introgressed
ancestry is largest for larger effect sizes (Figure~\ref{fig:linkage}C), and LD
between the selected and linked neutral alleles is largest for neutral sites
closest to the selected allele (Figure~\ref{fig:supp-LD}).

We assessed the accuracy of the deterministic two-locus model using
individual-based simulations \cite{thornton2019polygenic} under a simple
demographic model (Figure~\ref{fig:toy-admixture}A), with introgression
fraction \(f=0.05\). We simulated a single chromosome of length 1 Morgan, with
all mutations having effect sizes \(\pm a\) (\(a=0.02\) or \(0.05\), in
separate simulations), and we varied the total per-chromosome mutation rate
(\(\mu=0.001\), \(0.0025\) and \(0.01\)). For each fixed-difference mutation in
the parental populations, we determined the average introgressed ancestry
surrounding such loci $4{,}000$ generations after admixture (Methods).

The deterministic model (Equations~\ref{eq:system-p}--\ref{eq:system-D})
provides a very good approximation when mutation rates are small
(Figure~\ref{fig:linkage}D,E). As the mutation rate increases in these
simulations, trait-affecting alleles are more densely distributed along the
chromosome. Deviations from the deterministic model are due to multiple
selected alleles affecting local introgressed ancestry. For the highest
mutation rate shown here, there can be many other trait-affecting alleles
within a 1 cM window around any focal SNP, distorting dynamics away from
predictions under the two-locus model.

\subsection*{Introgressed ancestry deserts are shared under stabilizing
selection and reciprocal introgression}

\begin{figure}[ht!]
    \centering
    \caption{
        \textbf{Stabilizing selection causes an enrichment of introgressed
        ancestry deserts in functional regions.} (A) In chromosome-scale
        simulations under a demography with reciprocal migration between humans
        and Neanderthals (Figure~\ref{fig:human-neand-h2}A), stabilizing
        selection causes a chromosome-wide reduction of introgressed ancestry
        (below the $5\%$ and $2\%$ introgression proportions). This depletion
        is most pronounced in ``functional'' regions that allowed for
        trait-affecting mutations. (B) Introgressed ancestry deserts are
        more likely to occur in such functional regions, as are
        \emph{shared} deserts when compared across samples from humans and
        Neanderthals. In these simulations, $\sigma_M=0.02$ and $\mu=0.01$.
        Simulations with different mutational variances are shown in
        Figures~\ref{fig:deserts-high-VM}--\ref{fig:deserts-low-VM}.
        The pattern of shared deserts is not seen in simulations with
        directional selection (Figures~\ref{fig:deserts-directional-A}
        and \ref{fig:deserts-directional-B}).
    }
    \label{fig:deserts}
\end{figure}

As shown in the previous section, introgressed ancestry is reduced around loci
with trait-affecting alleles regardless of the parental population the derived
allele is present in. We should therefore expect to observe reductions in
introgressed ancestry at the same trait-affecting loci after gene flow in
either direction. If introgression occurs in both directions, regions of
reduced introgressed ancestry will coincide around such loci, and introgression
``deserts'' that appear due to selection against trait-affecting alleles will
tend to be shared.

To demonstrate this effect, we performed chromosome-scale simulations under a
model of reciprocal introgression between humans and Neanderthals
(Figure~\ref{fig:human-neand-h2}A and \cite{harris2023diverse}). Mutations
affecting a trait under stabilizing selection in each population occurred in
100 kb ``functional'' regions, spaced 2 Mb apart (Methods). Sampling
individuals from both the human and Neanderthal lineages, we observed the
average introgressed ancestry proportions were lowest within and surrounding
functional regions (Figure~\ref{fig:deserts}A). This corresponded to an
increased proportion of ancestry deserts within the functional regions
(Figure~\ref{fig:deserts}B). Functional regions also displayed an enrichment
of \emph{shared} ancestry deserts, with the probability of observing ancestry
deserts (either within a population or shared) decaying to background levels as
the distance from the functional region increases.

\section*{Discussion}

Many phenotypic traits are under stabilizing selection \cite{hodgins2015gene,
sanjak2018evidence, sella2019thinking}. This has motivated using stabilizing
selection around a shared optimum as a null model for the dynamics of alleles
affecting polygenic traits, including in multi-population settings
\cite{yair2022population}. Stabilizing selection on a trait results in
selection against the minor allele, mirroring symmetric underdominance
\cite{robertson1956effect, keightley1988quantitative}. Some theoretical and
simulation studies have considered the effect of population differentiation or
migration on the genetic architecture of a trait under stabilizing selection
(e.g., \cite{tufto2000quantitative, yeaman2011genetic, yair2022population}),
but most previous work has focused on single population scenarios, often
assuming steady state dynamics. Episodes of admixture and introgression
commonly occur in many species' evolutionary histories, so understanding their
effects on the architecture of complex traits is needed.

We find a multi-population, non-equilibrium approach for the site-frequency
spectrum with underdominance can be used to model the additive genetic variance
of a trait under stabilizing selection. While this approach ignores some
biological relevancies, such as linkage between sites, non-additivity and
pleiotropy, it still provides important insights into the dynamics of trait
architectures. For example, we can readily decompose contributions from alleles
of different origins, either by source population or mutation time. It may
therefore be useful for understanding the contributions of introgressed
variants to complex traits, such as hominin-introgressed mutations in humans
\cite{reilly2022contribution, wei2023lingering}. In the limited scenarios
explored here, we show that the expected contributions of introgressed variants
to heritability are complicated even in the purely additive case, depending on
the distribution of effect sizes, demographic history and the time since
admixture (Figures~\ref{fig:human-neand-h2} and \ref{fig:human-neand-VG}).

After admixture between diverged populations, the additive genetic variance of
a trait rapidly increases. The genetic variance after admixture depends on both
the existing genetic variances within the parental population and their
divergence at trait-affecting loci, measured by $F_2$
(Equation~\ref{eq:VG-admixture}) (see also, \cite{veller2024stabilizing}).
$V_G$ is then expected to decay fairly quickly to background levels. During
this process, introgressed and pre-existing trait-affecting alleles are
replaced by new mutations, turning over the genetic architecture of the trait
(Figure~\ref{fig:toy-admixture}). Selection occurs against the minor alleles at
trait-affecting loci, so that introgressed alleles from the minor parental
ancestry, whether derived or reintroduced ancestral alleles, are selected
against (Figure~\ref{fig:linkage}).

Selection against introgressed alleles also removes linked introgressed
ancestry in the surrounding regions. Thus, introgressed ancestry deserts are
more likely to form at and around loci contributing to selected traits
(Figure~\ref{fig:deserts}). This process is symmetric, so that deserts tend to
form in the same regions under reciprocal introgression. When comparing
distributions of introgressed haplotypes in two diverged populations with
recent bi-directional gene flow, we should expect to see an overlap of deserts
in genomic regions contributing to traits under stabilizing selection.

The expected pattern of shared ancestry deserts under stabilizing selection
differs from models of both deleterious load and incompatibilities, providing
testable hypotheses. Load-based models predict that haplotypes that have
accumulated more deleterious mutations, e.g., from a population with small
long-term effective population size, will be selected against under either
direction of gene flow. Introgressed ancestry at a given selected locus will
decrease in frequency in one introgression scenario and increase in the other.
This may explain the replacement of MT and Y chromosome DNA in Neanderthals by
human haplotypes after early human-to-Neanderthal introgression and the absence
of such Neanderthal haplotypes in modern humans \cite{posth2017deeply,
petr2020evolutionary}, but does not broadly match observations across the
autosomal genome \cite{harris2023diverse}.

The classical model of Bateson-Dobzhansky-Muller incompatibilities (BDMIs)
\cite{bateson1909heredity, dobzhansky1936studies, muller1942isolating}
explains the accumulation of reproductive isolation through negative epistatic
interactions that are exposed in hybrids. Muller \cite{muller1942isolating}
hypothesized that such BDMIs should most often form between distant or unlinked
loci, instead of within a single locus or tightly linked loci. Because theory
and experiments show that hybrid incompatibilities are resolved via selection
against the minor parental ancestry \cite{matute2020rapid, moran2021genomic},
ancestry deserts should form in different genomic regions, since different
incompatibility alleles are selected against depending on the parental ancestry
proportions. While both BDMI and stabilizing selection models predict selection
against introgressed alleles, an important distinction is that deserts due to
selection against pairs or small numbers of incompatibility loci are not
expected to overlap under bidirectional gene flow. While there is little
empirical data on the distribution of BDMI loci, studies point to interacting
BDMI alleles being unlinked \cite{presgraves2003fine, li2022imbalanced} and an
asymmetry in the alleles under selection in different introgression scenarios
\cite{maheshwari2011genetics, moran2021genomic}.

Recently, Harris \emph{et al.} \cite{harris2023diverse} observed that regions
of depleted Neanderthal ancestry in humans overlap more than would be expected
by chance with regions lacking human-introgressed alleles in the Altai
Neanderthal.  Human-introgressed ancestry in the Neanderthal genome is also
depleted in functional regions \cite{harris2023diverse}, as is well-documented
in humans \cite{sankararaman2014genomic, sankararaman2016combined}.  Harris
\emph{et al.} \cite{harris2023diverse} propose that epistatic interactions
between introgressed alleles and the recipient backgrounds could drive these
patterns, which they interpret as evidence for the initial process of
speciation between humans and Neanderthals. However, the observation of shared
ancestry deserts does not match expectations under a classic model of BDMIs as
described above.  Instead, at least some of the pattern may be due to
stabilizing selection acting on complex traits, such as gene regulation.
Importantly, such overlapping ancestry deserts are expected even when a trait
is under stabilizing selection for the same phenotypic optimal value. While the
underlying causes of selection on introgressed alleles in humans and
Neanderthals remain largely unknown, stabilizing selection provides a
well-grounded explanation for observed patterns that should be considered when
testing for epistasis, incompatibilities and adaptive introgression.

\section*{Methods}

\subsection*{Computing expected genetic variance from the SFS}

We model allelic dynamics using the diffusion approximation, where the expected
change in mean allele frequency per generation is
\[\E[M_{\delta_p}] = E[\Delta p] \approx s p(1-p),\]
with $s$ defined above, and the expected change in variance of the allele frequency per
generation is
\[\E[V_{\delta_p}] \approx \frac{1}{2N_e}p(1-p).\]


We extended the \moments-based solution for the sample site-frequency spectrum
(SFS) \cite{jouganous2017inferring} to include underdominance with given
selection coefficient $s$ (SI, Section~S2). The
contribution of alleles with effect size $a$ to the total genetic variance
$V_G$ is found by computing the expected pairwise diversity from the SFS (with
sample size $n$, denoted $\Phi_n$), as
\[V_{G,a} \approx 2a^2\sum_{j=1}^{n-1} \frac{j(n-j)}{n(n-1)} \Phi_n(j|a,\mu_a).\]
Here, $\mu_a$ is the mutation rate of alleles with effect size $a$. Then
assuming a normal distribution of effect sizes for new mutations, the total
genetic variance is
\[V_G \approx \int_{-\infty}^\infty V_{G,a} dF(a),\]
where \(F(a)\) is the cumulative distribution function of \(\mathcal{N}(0,V_M)\).
$V_G$ can be computed periodically in this way to predict trajectories in
non-equilibrium settings (Figure~\ref{fig:one-pop}).

If $V_G$ is non-negligible compared to $V_S$, ignoring $V_G$ and using
$s=(p-1/2)a^2/V_S$ leads to deflated estimates of additive genetic variance
(Figure~\ref{fig:supp-one-pop}). When mutation rates are large so that $V_G$ is
not small compared to $V_S$, using $s=(p-1/2)a^2/(V_S+V_G)$ provides estimates
of $V_G$ that closely match simulations assuming linkage equilibrium between
loci (Figures~\ref{fig:toy-admixture} and
\ref{fig:one-pop-0.01}--\ref{fig:one-pop-0.1}). Because $V_G$ can vary over
time, this means that $s$ is no longer constant and can change due to factors
such as non-constant demography that increase or reduce $V_G$.

\subsubsection*{Demographic history}

Our numerical solution for the SFS allows for non-constant population size
histories, population splits, continuous migration and admixture. Here, we
consider relatively simple scenarios involving population splits with
subsequent introgression events. We focus on parameter regimes relevant to
human-Neanderthal history. In a simple toy model a population of size
$N_e=10{,}000$ splits into two, one remaining size $10{,}000$ and the other
shrinking to size $1{,}000$. They remain isolated for $2N_e$ generations (or
500 thousand years, assuming an average generation time of 25 years) and then
introgression occurs from one branch to the other, contributing $5\%$ ancestry
to the recipient population (Figure~\ref{fig:toy-admixture}A,B).

The second model is meant to more closely resemble inferred human-Neanderthal
history, in which the ancestral population of size $N_e=10{,}000$ splits at 600
kya into the human and Neanderthal branches, with effective sizes $10{,}000$
and $2{,}000$, respectively. At 250 ka, an early human-to-Neanderthal
introgression contributes $5\%$ ancestry to Neanderthals. The human branch
shrinks to size $1{,}000$ 60 ka, followed by exponential growth to size
$20,000$ at present time. While present-day population sizes are much larger
\cite{tennessen2012evolution}, individual-based forward-in-time simulations
become computationally infeasible with increasing sizes, and this model should
still qualitatively capture evolutionary dynamics. Neanderthals contribute
$2\%$ ancestry to this bottlenecked and expanding human population at 50 ka,
after which they go extinct (Figure~\ref{fig:human-neand-h2}A).

In each scenario, we tracked phenotypic variance and genetic variation to
compare simulations to model predictions. The trait optimum was kept at $0$ in
all populations (no optimum shift occurred), and the strength of selection
$V_S=1$ remains constant.


\subsubsection*{Simulations with free recombination}

We compared our \moments-based predictions for $V_G$ in non-equilibrium
settings to simulations assuming linkage equilibrium as well as
individual-based simulations with recombination (next section). For both,
mutations occur at rate $\mu$ per haploid genome copy per generation, with
effects drawn from a normal distribution with mean 0 and variance $V_M$.

To simulate allele frequency changes under free recombination (assuming linkage
equilibrium between all trait-affecting alleles at all times), for each focal
segregating allele we integrated over possible genetic backgrounds contributed
by all other segregating alleles. Making the assumption that many alleles
contribute to the trait, the variance in genetic backgrounds is normally
distributed around the mean genetic value $E_{\tilde{G}}=\sum_l 2p_l a_l$ with
variance $V_{\tilde{G}}=\sum_l 2p_l(1-p_l)a_l^2$, where the sums omit the focal
locus. The expected change in frequency was then computed using the approach
outlined in SI, Section~S1 (without taking the
first-order Taylor series approximation for the exponentials). Allele counts in
the next generation were then binomially sampled with parameter $p+\E[\Delta
p]$ independently for each allele.

\subsubsection*{Individual-based simulations with linkage}

To include the effects of linkage between multiple selected alleles or selected
and neutral alleles, we used \fwdpy \cite{thornton2019polygenic} to run
Wright-Fisher simulations under the demographic models described above. In
these simulations, we considered large (1 Morgan, or 100 Mb with a per-base
recombination rate of $10^{-8}$) chromosomes with a uniform recombination
landscape. Trait-affecting mutations were either uniformly distributed across
the chromosome (Figures~\ref{fig:linkage} and
\ref{fig:supp-low-VM}--\ref{fig:supp-high-VM}) or fell within functional
regions (Figures~\ref{fig:deserts} and
\ref{fig:deserts-high-VM}--\ref{fig:deserts-directional-B}). For the latter,
such regions were centered 2 Mb apart and were 100kb in size, so that there
were 50 evenly spaced regions across the chromosome. Mutation effect sizes were
drawn either as constant values $\pm a$, or from a normal distribution with
mutational variance $V_M$.

In these simulations, we tracked the empirical phenotypic variance (since there
is no simulated environmental effect, this is equivalent to $V_G$) in each
population each generation. We measured the effects of linked selection on
neutral introgressed ancestry using \tskit to analyze genealogical
information \cite{ralph2020efficiently}. To measure the reduction in
introgressed ancestry around introgressed variants (as shown in
Figure~\ref{fig:linkage}), we identified each locus with a fixed difference
between the parental populations at the time of admixture by preserving the
generation immediately preceding admixture. For such a locus, for each sample
we determined which (preserved) parental population its ancestry traced to at
varying distances from the selected locus. Ancestry proportions were then
averaged over each fixed difference. To measure proportions of introgressed
ancestry or probabilities of observed ancestry deserts (as in
Figure~\ref{fig:deserts}), we again used \tskit to average ancestry proportions
of the source population in 50 kb windows. Ancestry deserts were defined as
any such window with no ancestry inherited from the source population.

\subsection*{Data and code availability}

All code to run analyses, create figures, and compile this manuscript is
available at
\url{https://github.com/apragsdale/neanderthal_stabilizing_selection}.


\section*{Acknowledgments}

I am grateful to Kevin Thornton, Bret Payseur, Carl Veller and Yuval Simons for
valuable discussions, and Carl Veller for helpful comments on an earlier
version of this manuscript. This work was supported by the National Institute
of General Medical Sciences of the National Institutes of Health (1R35GM154962
to A.P.R.). Support for this research was provided by the Office of the Vice
Chancellor for Research and Graduate Education at the University of
Wisconsin–Madison with funding from the Wisconsin Alumni Research Foundation.

\bibliographystyle{plos2015}
%\bibliography{manuscript}

\begin{thebibliography}{10}

\bibitem{brandvain2014speciation}
Brandvain Y, Kenney AM, Flagel L, Coop G, Sweigart AL.
\newblock {Speciation and introgression between \emph{Mimulus nasutus} and
  \emph{Mimulus guttatus}}.
\newblock PLoS genetics. 2014;10(6):e1004410.

\bibitem{skoglund2015ancient}
Skoglund P, Ersmark E, Palkopoulou E, Dal{\'e}n L.
\newblock {Ancient wolf genome reveals an early divergence of domestic dog
  ancestors and admixture into high-latitude breeds}.
\newblock Current Biology. 2015;25(11):1515--1519.

\bibitem{suvorov2022widespread}
Suvorov A, Kim BY, Wang J, Armstrong EE, Peede D, D’agostino ER, et~al.
\newblock {Widespread introgression across a phylogeny of 155 \emph{Drosophila}
  genomes}.
\newblock Current Biology. 2022;32(1):111--123.

\bibitem{tung2017contribution}
Tung J, Barreiro LB.
\newblock {The contribution of admixture to primate evolution}.
\newblock Current opinion in genetics \& development. 2017;47:61--68.

\bibitem{sorensen2023genome}
S{\o}rensen EF, Harris RA, Zhang L, Raveendran M, Kuderna LF, Walker JA, et~al.
\newblock {Genome-wide coancestry reveals details of ancient and recent
  male-driven reticulation in baboons}.
\newblock Science. 2023;380(6648):eabn8153.

\bibitem{wolf2018outstanding}
Wolf AB, Akey JM.
\newblock {Outstanding questions in the study of archaic hominin admixture}.
\newblock PLoS genetics. 2018;14(5):e1007349.

\bibitem{peter2020100}
Peter BM.
\newblock {100,000 years of gene flow between Neandertals and Denisovans in the
  Altai mountains}.
\newblock bioRxiv. 2020; p. 2020.03.13.990523.

\bibitem{prufer2014complete}
Pr{\"u}fer K, Racimo F, Patterson N, Jay F, Sankararaman S, Sawyer S, et~al.
\newblock {The complete genome sequence of a Neanderthal from the Altai
  Mountains}.
\newblock Nature. 2014;505(7481):43--49.

\bibitem{villanea2019multiple}
Villanea FA, Schraiber JG.
\newblock {Multiple episodes of interbreeding between Neanderthal and modern
  humans}.
\newblock Nature ecology \& evolution. 2019;3(1):39--44.

\bibitem{chen2020identifying}
Chen L, Wolf AB, Fu W, Li L, Akey JM.
\newblock {Identifying and interpreting apparent Neanderthal ancestry in
  African individuals}.
\newblock Cell. 2020;180(4):677--687.

\bibitem{sankararaman2016combined}
Sankararaman S, Mallick S, Patterson N, Reich D.
\newblock {The combined landscape of Denisovan and Neanderthal ancestry in
  present-day humans}.
\newblock Current Biology. 2016;26(9):1241--1247.

\bibitem{wei2023lingering}
Wei X, Robles CR, Pazokitoroudi A, Ganna A, Gusev A, Durvasula A, et~al.
\newblock {The lingering effects of Neanderthal introgression on human complex
  traits}.
\newblock Elife. 2023;12:e80757.

\bibitem{huerta2014altitude}
Huerta-S{\'a}nchez E, Jin X, Bianba Z, Peter BM, Vinckenbosch N, Liang Y,
  et~al.
\newblock {Altitude adaptation in Tibetans caused by introgression of
  Denisovan-like DNA}.
\newblock Nature. 2014;512(7513):194--197.

\bibitem{racimo2017signatures}
Racimo F, Marnetto D, Huerta-S{\'a}nchez E.
\newblock {Signatures of archaic adaptive introgression in present-day human
  populations}.
\newblock Molecular biology and evolution. 2017;34(2):296--317.

\bibitem{enard2018evidence}
Enard D, Petrov DA.
\newblock {Evidence that RNA viruses drove adaptive introgression between
  Neanderthals and modern humans}.
\newblock Cell. 2018;175(2):360--371.

\bibitem{gower2021detecting}
Gower G, Picazo PI, Fumagalli M, Racimo F.
\newblock {Detecting adaptive introgression in human evolution using
  convolutional neural networks}.
\newblock Elife. 2021;10:e64669.

\bibitem{harris2016genetic}
Harris K, Nielsen R.
\newblock {The genetic cost of Neanderthal introgression}.
\newblock Genetics. 2016;203(2):881--891.

\bibitem{juric2016strength}
Juric I, Aeschbacher S, Coop G.
\newblock {The strength of selection against Neanderthal introgression}.
\newblock PLoS genetics. 2016;12(11):e1006340.

\bibitem{veller2023recombination}
Veller C, Edelman NB, Muralidhar P, Nowak MA.
\newblock {Recombination and selection against introgressed DNA}.
\newblock Evolution. 2023;77(4):1131--1144.

\bibitem{kimura1963mutation}
Kimura M, Maruyama T, Crow JF.
\newblock The mutation load in small populations.
\newblock Genetics. 1963;48(10):1303.

\bibitem{lynch1990mutation}
Lynch M, Gabriel W.
\newblock Mutation load and the survival of small populations.
\newblock Evolution. 1990;44(7):1725--1737.

\bibitem{petr2019limits}
Petr M, P{\"a}{\"a}bo S, Kelso J, Vernot B.
\newblock {Limits of long-term selection against Neandertal introgression}.
\newblock Proceedings of the National Academy of Sciences.
  2019;116(5):1639--1644.

\bibitem{telis2020selection}
Telis N, Aguilar R, Harris K.
\newblock {Selection against archaic hominin genetic variation in regulatory
  regions}.
\newblock Nature Ecology \& Evolution. 2020;4(11):1558--1566.

\bibitem{yermakovich2023long}
Yermakovich D, Pankratov V, V{\~o}sa U, Yunusbayev B, {Estonian Biobank
  Research Team Metspalu Andres Nelis Mari Milani Lili M{\"a}gi Reedik Esko
  T{\~o}nu}, Dannemann M.
\newblock {Long-range regulatory effects of Neandertal DNA in modern humans}.
\newblock Genetics. 2023;223(3):iyac188.

\bibitem{sankararaman2014genomic}
Sankararaman S, Mallick S, Dannemann M, Pr{\"u}fer K, Kelso J, P{\"a}{\"a}bo S,
  et~al.
\newblock {The genomic landscape of Neanderthal ancestry in present-day
  humans}.
\newblock Nature. 2014;507(7492):354--357.

\bibitem{vernot2014resurrecting}
Vernot B, Akey JM.
\newblock Resurrecting surviving Neandertal lineages from modern human genomes.
\newblock Science. 2014;343(6174):1017--1021.

\bibitem{vernot2016excavating}
Vernot B, Tucci S, Kelso J, Schraiber JG, Wolf AB, Gittelman RM, et~al.
\newblock Excavating Neandertal and Denisovan DNA from the genomes of
  Melanesian individuals.
\newblock Science. 2016;352(6282):235--239.

\bibitem{harris2023diverse}
Harris DN, Platt A, Hansen ME, Fan S, McQuillan MA, Nyambo T, et~al.
\newblock {Diverse African genomes reveal selection on ancient modern human
  introgressions in Neanderthals}.
\newblock Current Biology. 2023;33(22):4905--4916.

\bibitem{kuhlwilm2016ancient}
Kuhlwilm M, Gronau I, Hubisz MJ, De~Filippo C, Prado-Martinez J, Kircher M,
  et~al.
\newblock {Ancient gene flow from early modern humans into Eastern
  Neanderthals}.
\newblock Nature. 2016;530(7591):429--433.

\bibitem{hubisz2020mapping}
Hubisz MJ, Williams AL, Siepel A.
\newblock {Mapping gene flow between ancient hominins through demography-aware
  inference of the ancestral recombination graph}.
\newblock PLoS genetics. 2020;16(8):e1008895.

\bibitem{li2024recurrent}
Li L, Comi TJ, Bierman RF, Akey JM.
\newblock {Recurrent gene flow between Neanderthals and modern humans over the
  past 200,000 years}.
\newblock Science. 2024;385(6705):eadi1768.

\bibitem{schwarcz1988esr}
Schwarcz HP, Gr{\"u}n R, Vandermeersch B, Bar-Yosef O, Valladas H, Tchernov E.
\newblock {ESR dates for the hominid burial site of Qafzeh in Israel}.
\newblock Journal of Human Evolution. 1988;17(8):733--737.

\bibitem{grun2005u}
Gr{\"u}n R, Stringer C, McDermott F, Nathan R, Porat N, Robertson S, et~al.
\newblock {U-series and ESR analyses of bones and teeth relating to the human
  burials from Skhul}.
\newblock Journal of human evolution. 2005;49(3):316--334.

\bibitem{harvati2019apidima}
Harvati K, R{\"o}ding C, Bosman AM, Karakostis FA, Gr{\"u}n R, Stringer C,
  et~al.
\newblock {Apidima Cave fossils provide earliest evidence of \emph{Homo
  sapiens} in Eurasia}.
\newblock Nature. 2019;571(7766):500--504.

\bibitem{beyer2021climatic}
Beyer RM, Krapp M, Eriksson A, Manica A.
\newblock {Climatic windows for human migration out of Africa in the past
  300,000 years}.
\newblock Nature communications. 2021;12(1):4889.

\bibitem{posth2017deeply}
Posth C, Wi{\ss}ing C, Kitagawa K, Pagani L, van Holstein L, Racimo F, et~al.
\newblock {Deeply divergent archaic mitochondrial genome provides lower time
  boundary for African gene flow into Neanderthals}.
\newblock Nature communications. 2017;8(1):16046.

\bibitem{petr2020evolutionary}
Petr M, Hajdinjak M, Fu Q, Essel E, Rougier H, Crevecoeur I, et~al.
\newblock {The evolutionary history of Neanderthal and Denisovan Y
  chromosomes}.
\newblock Science. 2020;369(6511):1653--1656.

\bibitem{muller1942isolating}
Muller HJ.
\newblock {Isolating mechanisms, evolution, and temperature}.
\newblock In: Biol Symp. vol.~6; 1942. p. 71--125.

\bibitem{sanjak2018evidence}
Sanjak JS, Sidorenko J, Robinson MR, Thornton KR, Visscher PM.
\newblock {Evidence of directional and stabilizing selection in contemporary
  humans}.
\newblock Proceedings of the National Academy of Sciences.
  2018;115(1):151--156.

\bibitem{sella2019thinking}
Sella G, Barton NH.
\newblock {Thinking about the evolution of complex traits in the era of
  genome-wide association studies}.
\newblock Annual review of genomics and human genetics. 2019;20:461--493.

\bibitem{gilad2006natural}
Gilad Y, Oshlack A, Rifkin SA.
\newblock {Natural selection on gene expression}.
\newblock Trends in Genetics. 2006;22(8):456--461.

\bibitem{hodgins2015gene}
Hodgins-Davis A, Rice DP, Townsend JP.
\newblock {Gene expression evolves under a house-of-cards model of stabilizing
  selection}.
\newblock Molecular Biology and Evolution. 2015;32(8):2130--2140.

\bibitem{price2022detecting}
Price PD, Palmer~Droguett DH, Taylor JA, Kim DW, Place ES, Rogers TF, et~al.
\newblock {Detecting signatures of selection on gene expression}.
\newblock Nature Ecology \& Evolution. 2022;6(7):1035--1045.

\bibitem{robertson1956effect}
Robertson A.
\newblock {The effect of selection against extreme deviants based on deviation
  or on homozygosis}.
\newblock Journal of Genetics. 1956;54(2):236--248.

\bibitem{keightley1988quantitative}
Keightley PD, Hill WG.
\newblock {Quantitative genetic variability maintained by mutation-stabilizing
  selection balance in finite populations}.
\newblock Genetics Research. 1988;52(1):33--43.

\bibitem{simons2018population}
Simons YB, Bullaughey K, Hudson RR, Sella G.
\newblock {A population genetic interpretation of GWAS findings for human
  quantitative traits}.
\newblock PLoS biology. 2018;16(3):e2002985.

\bibitem{hayward2022polygenic}
Hayward LK, Sella G.
\newblock {Polygenic adaptation after a sudden change in environment}.
\newblock Elife. 2022;11:e66697.

\bibitem{burger1989much}
B{\"u}rger R, Wagner GP, Stettinger F.
\newblock {How much heritable variation can be maintained in finite populations
  by mutation--selection balance?}
\newblock Evolution. 1989;43(8):1748--1766.

\bibitem{jouganous2017inferring}
Jouganous J, Long W, Ragsdale AP, Gravel S.
\newblock {Inferring the joint demographic history of multiple populations:
  beyond the diffusion approximation}.
\newblock Genetics. 2017;206(3):1549--1567.

\bibitem{yair2022population}
Yair S, Coop G.
\newblock {Population differentiation of polygenic score predictions under
  stabilizing selection}.
\newblock Philosophical Transactions of the Royal Society B.
  2022;377(1852):20200416.

\bibitem{veller2024stabilizing}
Veller C, Simons Y.
\newblock {Stabilizing selection generates selection against introgressed DNA}.
\newblock bioRxiv. 2024;.

\bibitem{castellano2014patterns}
Castellano S, Parra G, S{\'a}nchez-Quinto FA, Racimo F, Kuhlwilm M, Kircher M,
  et~al.
\newblock {Patterns of coding variation in the complete exomes of three
  Neandertals}.
\newblock Proceedings of the National Academy of Sciences.
  2014;111(18):6666--6671.

\bibitem{barton2002understanding}
Barton NH, Keightley PD.
\newblock {Understanding quantitative genetic variation}.
\newblock Nature Reviews Genetics. 2002;3(1):11--21.

\bibitem{zhang2002pleiotropic}
Zhang XS, Wang J, Hill WG.
\newblock {Pleiotropic model of maintenance of quantitative genetic variation
  at mutation-selection balance}.
\newblock Genetics. 2002;161(1):419--433.

\bibitem{lande1975maintenance}
Lande R.
\newblock {The maintenance of genetic variability by mutation in a polygenic
  character with linked loci}.
\newblock Genetics Research. 1975;26(3):221--235.

\bibitem{turelli1994genetic}
Turelli M, Barton NH.
\newblock {Genetic and statistical analyses of strong selection on polygenic
  traits: what, me normal?}
\newblock Genetics. 1994;138(3):913--941.

\bibitem{urban2013asymmetric}
Urban MC, B{\"u}rger R, Bolnick DI.
\newblock {Asymmetric selection and the evolution of extraordinary defences}.
\newblock Nature communications. 2013;4(1):2085.

\bibitem{lande1976natural}
Lande R.
\newblock {Natural selection and random genetic drift in phenotypic evolution}.
\newblock Evolution. 1976;30(2):314--334.

\bibitem{walsh2018evolution}
Walsh B, Lynch M.
\newblock {Evolution and selection of quantitative traits}.
\newblock Oxford University Press; 2018.

\bibitem{bulmer1971effect}
Bulmer M.
\newblock {The effect of selection on genetic variability}.
\newblock The American Naturalist. 1971;105(943):201--211.

\bibitem{negm2024effect}
Negm S, Veller C.
\newblock {The effect of long-range linkage disequilibrium on allele-frequency
  dynamics under stabilizing selection}.
\newblock bioRxiv. 2024; p. 2024--06.

\bibitem{peter2016admixture}
Peter BM.
\newblock {Admixture, population structure, and \emph{F}-statistics}.
\newblock Genetics. 2016;202(4):1485--1501.

\bibitem{tufto2000quantitative}
Tufto J.
\newblock {Quantitative genetic models for the balance between migration and
  stabilizing selection}.
\newblock Genetics Research. 2000;76(3):285--293.

\bibitem{burger1994distribution}
B{\"u}rger R, Lande R.
\newblock {On the distribution of the mean and variance of a quantitative trait
  under mutation-selection-drift balance}.
\newblock Genetics. 1994;138(3):901--912.

\bibitem{thornton2019polygenic}
Thornton KR.
\newblock {Polygenic adaptation to an environmental shift: temporal dynamics of
  variation under Gaussian stabilizing selection and additive effects on a
  single trait}.
\newblock Genetics. 2019;213(4):1513--1530.

\bibitem{hill1966effect}
Hill WG, Robertson A.
\newblock {The effect of linkage on limits to artificial selection}.
\newblock Genetics Research. 1966;8(3):269--294.

\bibitem{yeaman2011genetic}
Yeaman S, Whitlock MC.
\newblock {The genetic architecture of adaptation under migration--selection
  balance}.
\newblock Evolution. 2011;65(7):1897--1911.

\bibitem{reilly2022contribution}
Reilly PF, Tjahjadi A, Miller SL, Akey JM, Tucci S.
\newblock {The contribution of Neanderthal introgression to modern human
  traits}.
\newblock Current Biology. 2022;32(18):R970--R983.

\bibitem{bateson1909heredity}
Bateson W.
\newblock {Heredity and variation in modern lights}.
\newblock In: Darwin and modern science. Cambridge University Press; 1909. p.
  85--101.

\bibitem{dobzhansky1936studies}
Dobzhansky T.
\newblock {Studies on hybrid sterility. II. Localization of sterility factors
  in \emph{Drosophila pseudoobscura} hybrids}.
\newblock Genetics. 1936;21(2):113.

\bibitem{matute2020rapid}
Matute DR, Comeault AA, Earley E, Serrato-Capuchina A, Peede D, Monroy-Eklund
  A, et~al.
\newblock {Rapid and predictable evolution of admixed populations between two
  \emph{Drosophila} species pairs}.
\newblock Genetics. 2020;214(1):211--230.

\bibitem{moran2021genomic}
Moran BM, Payne C, Langdon Q, Powell DL, Brandvain Y, Schumer M.
\newblock {The genomic consequences of hybridization}.
\newblock Elife. 2021;10:e69016.

\bibitem{presgraves2003fine}
Presgraves DC.
\newblock {A fine-scale genetic analysis of hybrid incompatibilities in
  \emph{Drosophila}}.
\newblock Genetics. 2003;163(3):955--972.

\bibitem{li2022imbalanced}
Li J, Schumer M, Bank C.
\newblock {Imbalanced segregation of recombinant haplotypes in hybrid
  populations reveals inter-and intrachromosomal Dobzhansky-Muller
  incompatibilities}.
\newblock PLoS genetics. 2022;18(3):e1010120.

\bibitem{maheshwari2011genetics}
Maheshwari S, Barbash DA.
\newblock {The genetics of hybrid incompatibilities}.
\newblock Annual review of genetics. 2011;45:331--355.

\bibitem{tennessen2012evolution}
Tennessen JA, Bigham AW, O’connor TD, Fu W, Kenny EE, Gravel S, et~al.
\newblock Evolution and functional impact of rare coding variation from deep
  sequencing of human exomes.
\newblock science. 2012;337(6090):64--69.

\bibitem{ralph2020efficiently}
Ralph P, Thornton K, Kelleher J.
\newblock {Efficiently summarizing relationships in large samples: a general
  duality between statistics of genealogies and genomes}.
\newblock Genetics. 2020;215(3):779--797.

\end{thebibliography}

\section*{Supporting Information}

\renewcommand{\thesection}{S\arabic{section}}
\renewcommand{\thetable}{S\arabic{table}}
\renewcommand{\thefigure}{S\arabic{figure}}

\setcounter{figure}{0}

\begin{figure}[ht!]
    \centering
    \caption{
        Single-population dynamics when $V_G$ is not small.
        Solid lines: expected $V_G$ assuming $s=a^2/2V_S$. Dashed lines:
        expected $V_G$ assuming $s=a^2/2(V_S+V_G)$. When mutation rates
        are large (here, $\mu=0.01$, $V_S=1$), $V_G$ is non-negligible
        compared to $V_S$. By having to account for $V_G$ in the
        translation of effect size to selection coefficient, it makes
        $s$ non-constant if $V_G$ changes over time (due to non-constant
        demography, for example). In this case, $s$ must be updated
        regularly, given the current state of the population.
        Here, the demographic history is a bottleneck followed by a
        recovery, as depicted in Figure~\ref{fig:one-pop}A.
    }
    \label{fig:supp-one-pop}
\end{figure}

\begin{figure}[ht!]
    \centering
    \caption{
        Simulations with no linkage and weak effects.
        Total mutation rate was $\mu=0.025$ and all mutations had effect
        sizes $a=\pm 0.01$. The demographic history (bottleneck and
        recovery) is depicted in Figure~\ref{fig:one-pop}A. The predicted
        trajectory of additive genetic variance using \moments was found
        assuming symmetric underdominant selection on trait-affecting
        alleles, with selection coefficient $s=a^2/2(V_S+V_G)$.
    }
    \label{fig:one-pop-0.01}
\end{figure}

\begin{figure}[ht!]
    \centering
    \caption{
        Simulations with no linkage and moderate effects.
        All parameters were consistent with Figure~\ref{fig:one-pop-0.01},
        but with effect sizes $a=\pm0.04$.
    }
    \label{fig:one-pop-0.04}
\end{figure}

\begin{figure}[ht!]
    \centering
    \caption{
        Simulations with no linkage and strong effects.
        All parameters were consistent with Figure~\ref{fig:one-pop-0.01},
        but with effect sizes $a=\pm0.1$. In each comparison (weak, moderate,
        and strong effect sizes), predictions from \moments closely match
        observed average additive genetic variance from simulations.
    }
    \label{fig:one-pop-0.1}
\end{figure}


\begin{figure}[ht!]
    \centering
    \caption{
        With low mutational variance ($V_M=0.0001$), individual-based simulations
        using \fwdpy closely match our predictions using \moments, when the
        mutation rate is small (Left). When the mutation rate is increased,
        polygenicity increases and unlinked expectations from \moments deviate
        from observed $V_G$ in individual-based simulations. In this case, with
        relatively small mutational variance, the observed $V_G$ is reduced in
        comparison, consistent with the Bulmber (1971) effect.
    }
    \label{fig:supp-low-VM}
\end{figure}

\begin{figure}[ht!]
    \centering
    \caption{
        With high mutational variance ($V_M=0.0025$), strong deviations are
        observed between individual-based simulations and expectations without
        linkage. When the mutational input is large, observed $V_G$ can be
        much larger than expectations without linkage, possibly consistent with
        widespread interference between alleles.
    }
    \label{fig:supp-high-VM}
\end{figure}

\begin{figure}[t!]
    \centering
    \caption{
        Genetic variance of a trait under stabilizing selection in a model with
        human (blue) and Neanderthal (orange) demography. We consider two
        mutational variances ($\sigma_M=0.05$ and $0.01$). The demographic model
        (Figure~\ref{fig:human-neand-h2}A) includes $5\%$ admixture from humans
        to Neanderthals 250 ka, and $2\%$ admixture from Neanderthals to humans
        50 ka.
    }
    \label{fig:human-neand-VG}
\end{figure}


\begin{figure}[ht!]
    \centering
    \caption{
        Introgressed ancestry and probability of introgression deserts in and
        around functional regions, from simulations with high mutational variance.
        The demographic model is shown in Figure~\ref{fig:human-neand-h2}A.
    }
    \label{fig:deserts-high-VM}
\end{figure}

\begin{figure}[ht!]
    \centering
    \caption{
        Introgressed ancestry and probability of introgression deserts in and
        around functional regions, from simulations with moderate mutational variance.
        The demographic model is shown in Figure~\ref{fig:human-neand-h2}A.
    }
    \label{fig:deserts-med-VM}
\end{figure}

\begin{figure}[ht!]
    \centering
    \caption{
        Introgressed ancestry and probability of introgression deserts in and
        around functional regions, from simulations with low mutational variance.
        The demographic model is shown in Figure~\ref{fig:human-neand-h2}A.
    }
    \label{fig:deserts-low-VM}
\end{figure}

\begin{figure}[ht!]
    \centering
    \caption{
        Introgressed ancestry and probability of introgression deserts in and
        around functional regions, from simulations with directional selection.
        Deleterious mutations were drawn from a gamma distribution with mean
        $-0.0002$ and shape $0.05$.
        The demographic model is shown in Figure~\ref{fig:human-neand-h2}A.
    }
    \label{fig:deserts-directional-A}
\end{figure}

\begin{figure}[ht!]
    \centering
    \caption{
        Introgressed ancestry and probability of introgression deserts in and
        around functional regions, from simulations with directional selection.
        Deleterious mutations were drawn from a gamma distribution with mean
        $-0.001$ and shape $0.05$.
        The demographic model is shown in Figure~\ref{fig:human-neand-h2}A.
    }
    \label{fig:deserts-directional-B}
\end{figure}


\begin{figure}[ht!]
    \centering
    \caption{
        Underdominance, like negative selection, reduces expected $F_2$ compared
        to neutral divergence.
    }
    \label{fig:underdominance-F2}
\end{figure}

\begin{figure}[ht!]
    \centering
    \caption{
        Linkage disequilibrium between the a trait-affecting and neutral allele,
        2000 generations after admixture. Inititally, the introgression proportion
        was $0.05$. Mutations with stronger effect sizes, and thus stronger selection
        against them, decrease in frequency more rapidly, leading to reduced LD
        as measured by $D=Cov(p,q)$ (Figure~\ref{fig:linkage}).
    }
    \label{fig:supp-LD}
\end{figure}

\begin{figure}[ht!]
    \centering
    \caption{
        Comparison of predicted (\moments) and simulated SFS with underdominant
        selection. In this comparison, $N_e=10^4$ and $s=-5\times10^{-5}$, so that
        the population-size scale selection coefficient $\gamma=-1$.
        Simulations were performed under a Wright-Fisher model without linkage.
    }
    \label{fig:underdominance-validation-small-s}
\end{figure}

\begin{figure}[ht!]
    \centering
    \caption{
        Comparison of predicted (\moments) and simulated SFS with underdominant
        selection. In this comparison, $N_e=5000$ and $s=-0.001$, so that
        the population-size scale selection coefficient $\gamma=-10$.
        Simulations were performed under a Wright-Fisher model without linkage.
    }
    \label{fig:underdominance-validation-large-s}
\end{figure}

\begin{figure}[ht!]
    \centering
    \caption{
        Comparison of predicted (\moments) and simulated SFS with \emph{over}dominant
        selection. In this comparison, $N_e=5000$ and $s=0.0005$, so that
        the population-size scale selection coefficient $\gamma=5$.
        Thus, with positive selection, heterozygotes are favored.
        Simulations were performed under a Wright-Fisher model without linkage.
    }
    \label{fig:overdominance-validation}
\end{figure}

\end{document}
